% Imported from LaTeX_and_pdfs/AGAT_dump.tex; see provenance.json.
\part{Set Theory}


\section{Axioms of Set Theory}

\href{https://youtu.be/AAJB9l-HAZs?si=_yI6uRj7Vg1t0BbP}{ Axioms of set Theory - Lec 02 - Frederic Schuller}, Schuller (2016) \cite{agat:Schu2016}

\subsection{1.4 The $\epsilon$-Relation}

Set theory built on the postulate that there is a fundamental relation (i.e., a predicate of two variables) called $\epsilon$.

There will be no \emph{definition} of what $\epsilon$ is, or of what a set is.

Instead: 9 axioms that speak of $\epsilon$ and sets.

\subsubsection*{Overview of axioms}

EE PURP IC F

\[
\begin{gathered}
\begin{aligned}
	& \left. \begin{aligned}
		& E \\
		& E \\
	\end{aligned} \right\rbrace \text{ basic existence axioms } \\
& \left. \begin{aligned}
	& P \\
	& U \\
	& R \\
	& P \\
\end{aligned} \right\rbrace \text{ construction axioms } \\
& \left. \begin{aligned}
	& I \\
	& C \\
\end{aligned} \right\rbrace \text{ further existence / construction } \\
& \left. \begin{aligned}
	& F \\ 
\end{aligned} \right\rbrace \text{ non-existence }
\end{aligned}
\end{gathered}
\]

Using the $\epsilon$-relation, we can immediately define 
\[
\notin(x,y) =: x \notin y
\]

